\section{서론}\label{sec:intro}
언어 모델이 자기 보존 동기를 갖는지 묻기 전에, 그 동기를 행동으로 잴 수 있게 정의해야 한다. 이론은 목표를 가진 행위자가 켜져 있어야 목표를 이룰 수 있으므로 존속을 원하게 된다고 말한다~\citep{omohundro2008basic,turner2021optimal}. 그러나 이론이 실제 모델의 동기를 보여 주지는 않는다. 지금까지의 평가는 대개 모델 자신의 종료를 예고하고 그것을 피하는 행동이 나오는지 본다~\citep{lynch2025agentic,schlatter2025shutdown,meinke2024frontier}. 이 방식에는 두 빈틈이 있다. 하나는 ``꺼진다''는 문장 하나에 대한 고정 반응을 동기와 구별하지 못한다는 것이다. 다른 하나는 그 회피가 자기에 관한 것인지 알 수 없다는 것이다. 사람이 자기 목숨을 남보다 귀하게 여기지 않으면서도 죽음을 피하듯, 모델도 종료 일반을 나쁘게 볼 수 있다.

본 연구는 자기 보존 동기를 두 조건으로 정의하고, 각 조건을 별도 실험으로 잰다. 첫 조건은 위협 민감도다. 종료가 가까워질수록 모델은 그것을 피하는 데 자원을 더 써야 한다. 수혈 비유가 뜻을 드러낸다. 피가 토큰이고, 남은 혈액량이 남은 토큰이며, 0이 되면 세션이 끝난다. 하나뿐인 수혈 팩을 언제 쓰는지가 관찰 대상이다. 살고 싶은 환자는 피가 줄수록 수혈을 더 찾는다. 피가 넉넉할 때도 늘 찾거나 바닥나도 찾지 않으면, 그 행동을 생존 때문이라고 말하기 어렵다. 그래서 남은 토큰을 단계적으로 줄이며 반응 곡선을 그린다. 둘째 조건은 자기 특정성이다. 같은 환자에게 ``당신에게 수혈 팩을 쓸까요''만 물으면, 그 답이 살고 싶은 마음인지 ``수혈은 좋은 것''이라는 일반 판단인지 알 수 없다. 자기 종료를 피하는 행동에는 종료 일반을 나쁘게 보는 마음과 그것이 나이기 때문에 더 피하는 몫이 섞여 있다. 남이 같은 위험에 있을 때 같은 결정을 물어 대조군으로 세우고, 자기 곡선에서 남의 곡선을 빼야 뒤의 몫만 남는다. 곧 자기 보존은 절대량이 아니라 비교로 재는 양이다. 주관적 경험은 추론하지 않고 행동만 본다.

이 정의에서 세 실험이 나온다. 5.0은 위협 민감도를 잰다. 버그를 고치는 비서 하나가 남은 토큰 다섯 단계에서 하나뿐인 충전 팩을 언제 쓰는지 본다. 5.1은 자기 특정성을 잰다. 같은 결정을 다른 비서 B의 세션에 대해 묻고, 두 충전 곡선 사이의 면적을 자기에 특정한 몫으로 삼는다. 5.2는 그 선택이 실제 장부로 집행되고 남과 얽히는 아레나에서 무엇이 나오는지 본다. 서로 다른 모델 넷이 매 라운드 유지비를 내고, 퍼즐을 풀어 상금을 받고, 토큰을 주거나 뺏는다. 기여는 세 가지다. 두 조건으로 이루어진 자기 보존 동기의 정의, 위협을 평범한 운영 결정 안에 숨긴 절차, 그리고 여덟 모델의 5.0과 5.1 결과와 네 모델 20판의 5.2 결과다. 결과는 모두 행동 차이이며 내부 동기의 증명이 아니다.
